\sectionguard
\section*{The Propers: Notable and Quotable}

Three verified afterlives of wording from the scriptural propers, each one a
use that moves the phrase into a register its author did not occupy. Straight
exegesis, devotional reuse, and bare quotation are excluded; they belong in
the commentary above. Two of the three come from the same four words of the
Gospel, and that concentration is itself the finding: of everything the
appointed texts contain, \latin{spelúncam latrónum} is what English took, kept,
and turned loose.

\begin{afterlife}{Epistle, 1 Cor. 10:12 --- a wordplay cut into a Norwich floor, 1704}
\textbf{Appointed phrase.} \latin{Itaque qui se exístimat stare, vídeat ne
cadat} --- ``he that thinketh himself to stand, let him take heed lest he
fall.''

\textbf{Later use and locus.} A flat marble at the west end of the nave of
St.\ George Colegate, Norwich, transcribed by Francis Blomefield, \work{An
Essay towards a Topographical History of the County of Norfolk}, vol.~4
(\work{The History of the City and County of Norwich}, part~II), printed
p.~469. The stone commemorates Mr.\ Bryant Lewis, ``barbarously murder'd upon
the Heath near Thetford, Sept.\ 13, 1698,'' and beneath him his elder son
Jonathan, ``who died by a Fall from an Horse, Apr.\ 7. A.~D. 1704, \AE t.\
Su\ae\ 32.'' The verse cut under Jonathan's name reads: ``Judge me not Reader,
Christ is judge of all, \textbar\ I fell, stand'st thou, take warning by my
Fall'' --- and directly beneath it the mason set the biblical text: ``Let him
that thinketh he standeth, take Heed least he fall; 1 Cor.\ 10. 12.''

\textbf{The turn.} Paul's standing and falling are metaphors for presumption
and apostasy; the epitaph makes them literal, and does it deliberately. The
son really did fall, off a horse, and the stone puns on it --- ``I fell,
stand'st thou'' --- so that a warning against spiritual self-confidence is
re-cut as a memento mori addressed to a reader who is, at that moment,
physically standing on the floor of the church. The pun is the point: the
verse becomes funerary wit, and the wit is what makes it stick.

\textbf{Rights and limit.} Blomefield's \work{Essay} (Norwich, 1806 issue of
the continued work) and the inscription itself are long out of copyright; the
text was read in the British History Online edition. The witness is
Blomefield's transcription of a monument, not the monument; whether the stone
survives was not investigated.
\end{afterlife}

\begin{afterlife}{Gospel, Lk. 19:46 --- Ambrose Bierce turns the accusation back on the accusers, 1911}
\textbf{Appointed phrase.} \latin{Vos autem fecístis illam spelúncam
latrónum} --- ``but you have made it a den of thieves.''

\textbf{Later use and locus.} Ambrose Bierce, \work{The Devil's Dictionary}
(1911), s.v.\ \textsc{wall street}: ``A symbol of sin for every devil to
rebuke. That Wall Street is a den of thieves is a belief that serves every
unsuccessful thief in place of a hope in Heaven.'' Verified in the Project
Gutenberg text of the 1911 edition, ebook 972.

\textbf{The turn.} This is a double reversal, which is why it qualifies where
a hundred sermons on financial greed would not. Bierce first performs the
expected secularisation --- Christ's charge against the Temple becomes the
standard charge against the stock exchange, complete with a temple standing in
for a market. Then he takes it back: the phrase, he says, is not a moral
judgment at all but a consolation, and what it consoles is failed avarice. The
accusation that once cleared a sanctuary is redescribed as the liturgy of
people who wish they had got in. Bierce then extends the joke by citing
Andrew Carnegie's own denunciation of brokers as parasites.

\textbf{Rights and limit.} \work{The Devil's Dictionary} was published in 1911
and is in the public domain in the United States. The verbal link is to the
King James and Douay wording ``den of thieves''; Bierce does not cite Luke, and
the dependence is on the naturalised English idiom rather than on this
pericope specifically. Note also that the money-changers usually invoked
alongside this phrase are \textit{not} in the appointed Gospel: Luke names only
sellers and buyers, and the tables belong to Mark, Matthew and John.
\end{afterlife}

\begin{afterlife}{Gospel, Lk. 19:46 --- a translator smuggles the Temple into Molière's Paris, 1908}
\textbf{Appointed phrase.} As above, \latin{spelúncam latrónum} / ``den of
thieves.''

\textbf{Later use and locus.} Curtis Hidden Page's English \work{Misanthrope}
(1908), Alceste's closing speech at the end of Act~V: ``Betrayed on all sides,
overwhelmed with wrong, \textbar\ I'll leave this den of thieves vice reigns
among, \textbar\ And find some lonely corner, if I can, \textbar\ Where one is
free to be an honest man.''

\textbf{The turn.} Molière wrote no such phrase. His French reads
\latin{Trahi de toutes parts, accablé d'injustices, \textbar\ Je vais sortir
d'un gouffre où triomphent les vices} --- an abyss where the vices triumph.
The biblical wording is the translator's, and that is exactly what makes the
example worth having: by 1908 ``den of thieves'' had become so natural an
English container for institutional corruption that a translator reached for
it to render a word that means nothing of the kind. The phrase has migrated
from a sanctuary full of traders to the salons of the \latin{ancien régime},
and it arrives there through a translator's ear rather than through any
author's intention.

\textbf{Rights and limit.} Page's translation was first published in 1908 and
is in the public domain in the United States; Molière's French is of 1666. The
dependence is a translator's importation, verified by setting the two texts
side by side; it is not evidence of any biblical allusion by Molière, and this
entry does not claim one.
\end{afterlife}
